% from ega4-14-fr.tex
\oldpage[IV-3]{215}

un voisinage ouvert $U$ de $x$ dans $X$~; comme $v$ est fermé et surjectif, pour
que $f(U)$ soit un voisinage de $y=f(x)$, il faut et il suffit que
$v^{-1}(f(U))$ soit un voisinage de $v^{-1}(y)$. Mais si $p:X_1\to X$ est la
projection canonique, on a $v^{-1}(f(U))=f_1(p^{-1}(U))$ (I, 3.4.8), et
l'hypothèse entraîne que $f_1(p^{-1}(U))$ est un voisinage de
$f_1(p^{-1}(x))=v^{-1}(y)$ (I, 3.4.8).

\begin{corollary}[14.4.9]
\label{IV.14.4.9-fr}
Soient $Y$ un préschéma noethérien, $f:X\to Y$ un morphisme localement de type
fini. Les conditions suivantes sont équivalentes~:
\begin{enumerate}[label=\emph{\alph*)}]
  \item $f$ est universellement ouvert, autrement dit, pour tout changement de
  base $Y'\to Y$, le morphisme $f_{(Y')}:X\times_Y Y'\to Y'$ est ouvert.
  \item[a$'$)] Pour tout morphisme fini $Y_1\to Y$, $f_{(Y_1)}$ est ouvert.
  \item[a$''$)] Si $Y''$ est le normalisé de $Y_{\mathrm{red}}$, $f_{(Y'')}$ est
  ouvert.
  \item Pour tout point $x''$ de $X''=X\times_Y Y''$, il existe une composante
  irréductible $Z''$ de $X''$ contenant $x''$ et équidimensionnelle sur $Y''$ au
  point $x''$ (cf. (14.4.10, (ii))).
\end{enumerate}
\end{corollary}

Cela résulte aussitôt de (14.4.8) et (14.1.4).

\begin{remarks}[14.4.10]
\label{IV.14.4.10-fr}
\begin{enumerate}[label=(\roman*)]
  \item L'équivalence des conditions $a)$ et $b)$ dans (14.4.8) (resp. (14.4.9))
  est encore valable pour un préschéma $Y$ quelconque et un morphisme $f$
  localement de type fini. En effet, $a)$ entraîne $b)$ en vertu de (14.4.1)~;
  inversement, $b)$ entraîne que $f''$ est universellement ouvert aux points de
  $X''$ au-dessus de $x$ (resp. en tout point de $X''$) en vertu de (14.4.1), et
  on en conclut la propriété $a)$ en appliquant le lemme (14.4.8.1) au morphisme
  entier et surjectif $Y''\to Y$.

  Il se peut que, dans (14.4.1), pour l'équivalence de $a)$ et $c)$, l'hypothèse
  supplémentaire que $Y$ est noethérien soit superflue (cf. (14.3.14)). S'il en
  est ainsi, les hypothèses noethériennes sont également superflues dans (14.4.2),
  (14.4.3), (14.4.8) et (14.4.9).

  \item On peut donner des exemples de morphismes $f:X\to Y$ ayant les propriétés
  suivantes~: $Y$ est noethérien, régulier et de dimension 2, $f$ est
  \emph{universellement ouvert} et de type fini, $X$ a deux composantes
  irréductibles $X_1$, $X_2$, mais la restriction $X_1\to Y$ de $f$ à l'une
  d'elles \emph{n'est pas un morphisme ouvert}. Le principe de la construction
  s'appuie sur la méthode générale de «~recollement~» qui sera exposée au chap.~V,
  et ne peut donc qu'être esquissé ici. On part d'un point fermé $y$ de $Y$, et on
  considère le $Y$-schéma $Y_1$ obtenu en faisant \emph{éclater} $y$
  (\textbf{II}, 8.1.3)~; si $f_1:Y_1\to Y$ est le morphisme structural, on sait que
  la restriction de $f_1$ à $f_1^{-1}(Y-\{y\})$ est un \emph{isomorphisme} sur
  $Y-\{y\}$ (\emph{loc. cit.}), tandis que la fibre $f_1^{-1}(y)$ est isomorphe à
  $\operatorname{Proj}(S)$, où
  $S=\bigoplus_{k=0}^{\infty}\mathfrak m_y^k/\mathfrak m_y^{k+1}$
  (\textbf{II}, 3.5.3), c'est-à-dire ici à $\mathbf P^1_{k(y)}$~; il résulte de
  (14.4.1) que $f_1$ n'est pas ouvert au point générique de $f_1^{-1}(y)$. D'autre
  part, posons $Y_2=\mathbf P^1_Y$, et soit $f_2:Y_2\to Y$ le morphisme structural~;
  il résulte de (\textbf{II}, 8.4.4) que $f_2$ est plat, donc universellement ouvert
  (2.4.6)~; en outre (\textbf{II}, 3.5.3), $f_2^{-1}(y)$ est isomorphe à
  $\mathbf P^1_{k(y)}$~; il suffit alors de «~recoller~» $Y_1$ et $Y_2$ le long des
  fibres isomorphes $f_1^{-1}(y)$ et $f_2^{-1}(y)$, ce qui donne un morphisme
  $f:X\to Y$ où les composantes irréductibles $X_1$, $X_2$ de $X$ s'identifient
  canoniquement à $Y_1$ et $Y_2$ respectivement, et les restrictions de $f$ à ces
  composantes à $f_1$ et $f_2$.

  Rappelons toutefois (12.1.1.5) que si $Y$ est localement noethérien, $f:X\to Y$
  de type fini et \emph{plat}, alors toute composante irréductible de $X$ est
  \emph{équidimensionnelle} sur $Y$ (et par suite la restriction de $f$ à une telle
  composante est \emph{universellement ouverte} si tous les points de $Y$ sont
  \emph{géométriquement unibranches}).

  \item Rappelons (12.1.2, (i)) qu'il y a des morphismes $f:X\to Y$ ayant les
  propriétés suivantes~: $Y$ est noethérien (non géométriquement unibranche), $f$
  est fini et plat (et même étale (17.6.3)), mais la restriction de $f$ à une
  composante irréductible de $X$ \emph{n'est pas un morphisme ouvert} (bien que $f$
  lui-même le soit d'après (2.4.6)).

  \item Le critère de Chevalley (14.4.4) explique l'importance de la notion de
  morphisme universellement ouvert. Cette notion permet en effet, dans de nombreux
  résultats, plus ou moins classiques, de remplacer une hypothèse de \emph{normalité}
  par l'hypothèse qu'un certain morphisme est universellement ouvert~; l'énoncé plus
  général

  